\chapter{Pendahuluan}

\section{Latar Belakang}

Kebutuhan listrik global tumbuh sebesar 4,4\% pada tahun 2024 dan diproyeksikan
meningkat rata-rata 3,6\% per tahun sepanjang 2026--2030, sekitar 50\% lebih
tinggi dibandingkan rata-rata dekade sebelumnya, didorong oleh elektrifikasi
industri, kendaraan listrik, pendinginan udara, dan pusat data
\cite{iea2026electricity}. EBT dan nuklir bersama-sama diproyeksikan mencakup
sekitar setengah dari total pembangkitan listrik dunia pada tahun 2030, dengan
EBT diperkirakan memenuhi hampir seluruh tambahan kebutuhan listrik tersebut
\cite{iea2026electricity}. Di Indonesia, pemerintah menargetkan bauran EBT
sebesar 23\% pada tahun 2025 \cite{iea2022indonesia}. Rencana Usaha Penyediaan
Tenaga Listrik (RUPTL) PLN 2025--2034 menetapkan penambahan 69,5 GW kapasitas
pembangkit, 76\% di antaranya dari EBT \cite{ruptl2025}.

Secara global, kapasitas terpasang Pembangkit Listrik Tenaga Bayu (PLTB) telah
mencapai 1.133 GW pada akhir tahun 2024 dengan penambahan sebesar 113,2 GW
sepanjang tahun tersebut, menjadikan energi angin sebagai sumber EBT dengan
kapasitas terbesar kedua di dunia setelah energi surya \cite{irena2025capacity}.
Di Indonesia, potensi energi angin nasional diperkirakan mencapai 154,6 GW,
namun kapasitas PLTB yang telah terpasang hingga tahun 2024 baru sebesar 152,3
MW, atau kurang dari 0,1\% dari total potensi yang tersedia
\cite{esdm2025pltb,pambudi2025wind}. Kesenjangan ini bersumber dari harga
listrik PLTB yang lebih tinggi dibandingkan pembangkit konvensional,
keterbatasan industri komponen lokal, dan minimnya pembiayaan berbiaya rendah
\cite{pambudi2025wind}. Pemerintah Indonesia menetapkan target penambahan 5 GW
kapasitas PLTB hingga tahun 2030 \cite{esdm2025pltb}. Skala target ini menuntut
pemahaman yang lebih baik tentang dampak integrasi PLTB terhadap stabilitas
sistem tenaga listrik, termasuk \textit{small signal stability}.

Integrasi PLTB dalam skala besar menimbulkan tantangan teknis, khususnya terhadap
\textit{small signal stability} (SSS) atau stabilitas sinyal kecil. Integrasi
PLTB, baik berbasis \textit{doubly fed induction generator} (DFIG)/Type-3 maupun
\textit{full-rated converter} (FRC)/Type-4, mengubah karakteristik dinamik sistem,
namun arah perubahan tersebut bersifat \textit{system-dependent} dan tidak universal. Peningkatan penetrasi DFIG pada sistem berskala besar dapat memberikan dampak
ganda terhadap \textit{electromechanical oscillation mode}, yakni menguntungkan
pada sejumlah \textit{oscillation mode} namun merugikan pada
\textit{oscillation mode} lainnya, bergantung pada lokasi integrasi dan
sensitivitas \textit{eigenvalue} terhadap inersia sistem \cite{gautam2009impact}. Perbandingan tiga jenis \textit{wind turbine generator} (WTG) yaitu
\textit{squirrel cage induction generator} (SCIG), DFIG, dan \textit{permanent
magnet synchronous generator} (PMSG) pada sistem \textit{nine-bus} WECC dan
sistem 24-bus menunjukkan bahwa dampak SSS dari masing-masing teknologi
bergantung pada strategi kontrol dan kondisi operasi, di mana DFIG dan PMSG
berpotensi memunculkan \textit{eigenvalue} bernilai positif yang
mengindikasikan ketidakstabilan pada tingkat penetrasi tinggi
\cite{he2015investigation}. Kajian pada sistem WECC 9-bus termodifikasi menunjukkan bahwa peningkatan
keluaran DFIG memperbaiki SSS \cite{sun2013small}. Melalui pendekatan
\textit{stability region boundary}, dampak DFIG terhadap SSS terbukti bergantung
pada skala penetrasi dan kekuatan kopling dengan generator sinkron
\cite{liu2014impact}. Untuk teknologi PMSG-WT, turbin angin berbasis konverter penuh memiliki
partisipasi yang rendah terhadap osilasi antar-area karena konverter penuh
memisahkan dinamika generator dari dinamika jaringan \cite{knuppel2009small}.
Model SSS untuk turbin angin \textit{direct-drive} berbasis PMSG menunjukkan
bahwa teknologi ini menghasilkan respons dinamik yang lebih baik dibandingkan
DFIG \cite{wu2009small}. Kajian yang langsung membandingkan dampak DFIG-WT
dan PMSG-WT terhadap SSS pada sistem \textit{benchmark} yang sama, dengan
variasi posisi bus integrasi dan kondisi pembebanan, serta metrik meliputi
\textit{eigenvalue}, \textit{damping ratio}, dan \textit{participation factor},
masih terbatas.

Analisis SSS melinearisasi model sistem di sekitar titik operasi keseimbangan
dan menilai kestabilan terhadap gangguan kecil melalui nilai \textit{eigenvalue},
\textit{damping ratio}, dan \textit{participation factor}. Bagian real dari
\textit{eigenvalue} menentukan sifat redaman \textit{oscillation mode}: nilai
negatif menunjukkan \textit{oscillation mode} yang stabil, sedangkan nilai
positif mengindikasikan ketidakstabilan. Bagian
imajiner menentukan frekuensi \textit{electromechanical oscillation} yang bersesuaian.
\textit{Damping ratio} ($\zeta$) mengukur tingkat redaman relatif, dengan ambang batas
$\zeta \geq 5\%$ umumnya digunakan sebagai kriteria kestabilan yang memadai.
\textit{Participation factor} mengidentifikasi variabel keadaan yang paling
berpengaruh terhadap suatu \textit{oscillation mode}, sehingga memungkinkan pelacakan
kontribusi masing-masing generator terhadap dinamika sistem. Penerapan kerangka
ini pada sistem uji IEEE 9-bus dengan penetrasi tinggi PLTB berbasis PMSG
menunjukkan bahwa pendekatan modal berbasis \textit{eigenvalue} mampu
mengidentifikasi pergeseran \textit{oscillation mode} dominan pada berbagai tingkat penetrasi
\cite{wang2014small}.

Penelitian ini menganalisis dan membandingkan dampak integrasi PLTB berbasis
DFIG-WT dan PMSG-WT terhadap \textit{small signal stability} sistem tenaga
listrik pada sistem uji \textit{nine-bus} IEEE. Pemodelan dan simulasi
menggunakan perangkat lunak DIgSILENT PowerFactory, dengan analisis modal
berbasis \textit{eigenvalue} untuk mengidentifikasi \textit{oscillation mode} dominan,
\textit{damping ratio}, frekuensi osilasi, dan \textit{participation factor}
dari setiap skenario. Pengaruh posisi integrasi dikaji dengan menempatkan PLTB
pada dua lokasi yang berbeda, masing-masing menggantikan generator konvensional
yang ada, sementara kondisi pembebanan divariasikan untuk mencerminkan berbagai titik operasi.
Penelitian ini menghasilkan perbandingan kuantitatif dampak kedua teknologi WTG
terhadap SSS dan masukan teknis bagi perencanaan integrasi PLTB ke dalam sistem
tenaga listrik nasional.

\section{Rumusan Masalah}

Penelitian ini menjawab tiga pertanyaan berikut:

\begin{enumerate}
	\item Bagaimana pengaruh integrasi PLTB berbasis DFIG-WT terhadap
	\textit{small signal stability} sistem tenaga listrik pada variasi posisi bus
	integrasi dan kondisi pembebanan?
	\item Bagaimana pengaruh integrasi PLTB berbasis PMSG-WT terhadap
	\textit{small signal stability} sistem tenaga listrik pada variasi posisi bus
	integrasi dan kondisi pembebanan?
	\item Bagaimana perbandingan dampak integrasi DFIG-WT dan PMSG-WT terhadap
	\textit{small signal stability} ditinjau dari \textit{eigenvalue},
	\textit{damping ratio}, dan \textit{participation factor}?
\end{enumerate}

\section{Tujuan Penelitian}

Penelitian ini bertujuan untuk:

\begin{enumerate}
	\item Menganalisis dampak integrasi PLTB berbasis DFIG-WT terhadap
	\textit{electromechanical oscillation mode}, \textit{damping ratio}, dan
	\textit{participation factor} pada sistem uji \textit{nine-bus} IEEE.
	\item Menganalisis dampak integrasi PLTB berbasis PMSG-WT terhadap
	\textit{electromechanical oscillation mode}, \textit{damping ratio}, dan
	\textit{participation factor} pada sistem uji \textit{nine-bus} IEEE.
	\item Membandingkan secara kuantitatif dampak integrasi DFIG-WT dan PMSG-WT
	terhadap \textit{small signal stability} pada berbagai posisi bus integrasi
	dan kondisi pembebanan.
\end{enumerate}


\section{Batasan Penelitian}

Penelitian ini dibatasi pada:

\begin{enumerate}
	\item \textbf{Sistem Uji:} Sistem yang digunakan adalah sistem uji standar
	IEEE \textit{nine-bus} dengan parameter yang telah ditetapkan, beroperasi
	pada frekuensi nominal 60 Hz, tanpa modifikasi topologi jaringan.
	\item \textbf{Lingkup Analisis Kestabilan:} Analisis dibatasi pada
	\textit{small signal stability}; kestabilan transien (\textit{transient
	stability}) dan kestabilan tegangan (\textit{voltage stability}) tidak
	dikaji. Sistem diasumsikan beroperasi pada kondisi normal tanpa gangguan
	besar (\textit{large disturbance}).
	\item \textbf{Metode Analisis:} Analisis dilakukan secara eksklusif
	menggunakan analisis modal berbasis \textit{eigenvalue} (\textit{modal
	analysis}), tanpa simulasi RMS (\textit{root mean square}) berbasis
	domain waktu. Oleh karena itu, hasil penelitian ini tidak mencakup plot
	respons waktu (\textit{time response}) sistem terhadap gangguan.
	\item \textbf{Pemodelan PLTB:} Model komponen dan parameter kontroler
	DFIG-WT serta PMSG-WT seluruhnya menggunakan nilai bawaan
	(\textit{default}) dari pustaka DIgSILENT PowerFactory 2024, tanpa
	modifikasi atau tuning tambahan.
	\item \textbf{Skenario Integrasi:} PLTB hanya menggantikan satu generator
	sinkron konvensional (bukan \textit{slack bus}) dengan kapasitas yang
	setara; setiap skenario mengkaji satu tingkat penetrasi tanpa variasi
	penetrasi parsial.
	\item \textbf{Asumsi Pembebanan dan Angin:} Kecepatan angin diasumsikan
	konstan selama simulasi. Variasi beban diterapkan pada satu beban spesifik;
	beban lain tidak diubah.
	\item \textbf{Aspek yang Tidak Dimodelkan:} Pengaruh harmonik dari
	konverter \textit{power electronics}, serta \textit{Automatic Voltage
	Regulator} (AVR) dan \textit{Power System Stabilizer} (PSS) pada seluruh
	generator sinkron konvensional, tidak diperhitungkan dalam analisis.
	\item \textbf{Perangkat Lunak:} Penelitian menggunakan DIgSILENT
	PowerFactory 2024 sebagai satu-satunya alat simulasi.
\end{enumerate}


%\noindent \textbf{Contoh penulisan batasan skripsi Teknik Elektro:}
%
%%------------------------------------------------- 
%\noindent\fbox{% \parbox{\textwidth}{% \begin{enumerate} \item Objek
%   penelitian: Studi efisiensi dan kinerja sistem pemantauan suhu dan arus pada
%   sistem pembangkit tenaga listrik. \item Metode penelitian: Penelitian
%   eksperimental menggunakan analisis simulasi dan pengujian sistem pada skala
%   laboratorium. \item Waktu dan tempat penelitian: Waktu penelitian adalah
%   Maret-Agustus 2022 di FasLab TTL. \item Populasi dan sampel: Populasi adalah
%   sistem pembangkit tenaga listrik, dan sampel diambil sebanyak 3 sistem yang
%   berbeda. \item Variabel: Variabel bebas adalah metode pemantauan suhu dan
%   arus, dan variabel terikat adalah efisiensi dan kinerja \item Hipotesis:
%   bahwa metode pemantauan suhu dan arus berpengaruh terhadap efisiensi dan
%   kinerja sistem pembangkit tenaga listrik. \item Keterbatasan Penelitian:
%   Keterbatasan penelitian adalah hanya melibatkan pengujian sistem pada skala
%   laboratorium dan membatasi analisis pada variabel bebas dan terikat.
%   \end{enumerate}
%
%   }%
%}
%
%%------------------------------------------------- 
%\newpage
%
%\noindent \textbf{Contoh penulisan batasan skripsi Teknik Biomedis:}
%
%%------------------------------------------------- 
%\noindent\fbox{% \parbox{\textwidth}{% \begin{enumerate} \item Objek
%   Penelitian: Studi efektivitas dan keamanan alat elektromedik seperti pacu
%   jantung. \item Metode Penelitian: Penelitian deskriptif dan observasional
%   menggunakan survei dan analisis data. \item Waktu dan Tempat Penelitian:
%   Waktu penelitian adalah Januari-Juni 2022 di Rumah Sakit Sarjito. \item
%   Populasi dan Sampel: Populasi adalah pasien yang menggunakan pacu jantung,
%   dan sampel diambil dengan metode random sampling sebanyak 100 pasien. \item
%   Variabel: Variabel bebas nya adalah jenis alat pacu jantung, dan variabel
%   terikatnya adalah efektivitas dan keamanan. \item Hipotesis: bahwa jenis
%   alat pacu jantung berpengaruh terhadap efektivitas dan keamanan. \item
%   Keterbatasan Penelitian: Keterbatasan penelitian adalah hanya melibatkan
%   satu rumah sakit sebagai lokasi penelitian dan membatasi dalam analisis data
%   hanya pada variabel bebas dan terikat. \end{enumerate}
%
%   }%
%}
%
%%------------------------------------------------- 
%
%\newpage \noindent \textbf{Contoh penulisan batasan skripsi Teknologi
%Informasi:}
%
%%------------------------------------------------- 
%\noindent\fbox{% \parbox{\textwidth}{% \begin{enumerate} \item Objek
%   Penelitian: Analisis perbandingan efektivitas dan efisiensi antara sistem
%   manajemen proyek tradisional dan sistem manajemen proyek berbasis teknologi
%   informasi. \item Metode Penelitian: Penelitian kualitatif dengan menggunakan
%   wawancara dan survei terhadap para pelaku proyek di berbagai perusahaan.
%   \item Waktu dan Tempat Penelitian: Waktu penelitian adalah Januari-Juni 2022
%   di perusahaan-perusahaan di wilayah Bantul. \item Populasi dan Sampel:
%   Populasi nya adalah perusahaan yang melakukan proyek, dan sampel diambil
%   sebanyak 10 perusahaan yang menerapkan sistem manajemen proyek tradisional
%   dan 10 perusahaan yang menggunakan sistem manajemen proyek berbasis
%   teknologi informasi. \item Variabel: Variabel bebasnya adalah sistem
%   manajemen proyek, dan variabel terikatnya adalah efektivitas dan efisiensi.
%   \item Hipotesis: bahwa sistem manajemen proyek berbasis teknologi informasi
%   lebih efektif dan efisien dibandingkan dengan sistem manajemen proyek
%   tradisional. \item Keterbatasan Penelitian: Keterbatasan penelitian adalah
%   penelitian hanya dilakukan pada perusahaan di wilayah Bantul dan hanya
%   melibatkan wawancara dan survei sebagai metode pengumpulan data.
%   \end{enumerate}
%
%   }%
%}

\section{Manfaat Penelitian}

Penelitian ini memberi tiga manfaat:

\begin{enumerate}
	\item \textbf{Akademis/Teoritis:} Perbandingan kuantitatif DFIG-WT dan
	PMSG-WT terhadap \textit{small signal stability} pada IEEE \textit{nine-bus}
	menjadi acuan bagi peneliti kestabilan sistem tenaga dengan penetrasi EBT.
	\item \textbf{Analitis:} Analisis \textit{eigenvalue}, \textit{damping
	ratio}, dan \textit{participation factor} mengungkap variabel keadaan yang
	mendominasi \textit{oscillation mode} saat PLTB terhubung ke sistem.
	\item \textbf{Industri/Praktis:} Perencana sistem tenaga nasional mendapat
	acuan teknis untuk memilih teknologi WTG dan posisi integrasi PLTB menuju
	target 5 GW pada 2030.
\end{enumerate}



\section{Sistematika Penulisan}

\noindent \textbf{Bab I Pendahuluan} memuat latar belakang yang menjelaskan
urgensi integrasi PLTB terhadap stabilitas sistem tenaga listrik, rumusan masalah
yang merinci tiga pertanyaan penelitian, tujuan penelitian, batasan penelitian,
manfaat penelitian, serta sistematika penulisan laporan.

\noindent \textbf{Bab II Tinjauan Pustaka dan Dasar Teori} menguraikan kajian
literatur terkait dampak integrasi PLTB berbasis DFIG-WT dan PMSG-WT terhadap
\textit{small signal stability} (SSS), serta memaparkan dasar teori yang meliputi
pemodelan sistem tenaga listrik, karakteristik teknologi WTG, konsep analisis
modal, dan metrik kestabilan yang digunakan yaitu \textit{eigenvalue},
\textit{damping ratio}, dan \textit{participation factor}.

\noindent \textbf{Bab III Metode Penelitian} menjelaskan pendekatan penelitian
kuantitatif berbasis simulasi, alat dan bahan penelitian, pemodelan sistem uji
IEEE \textit{nine-bus} pada DIgSILENT PowerFactory 2024, konfigurasi model
DFIG-WT dan PMSG-WT, serta definisi enam skenario percobaan yang mencakup
variasi posisi bus integrasi dan variasi kondisi pembebanan.

\noindent \textbf{Bab IV Hasil dan Pembahasan} menyajikan dan menganalisis hasil
simulasi dari seluruh skenario percobaan, meliputi nilai \textit{eigenvalue},
\textit{damping ratio}, frekuensi osilasi, dan \textit{participation factor} pada
masing-masing skenario, serta membahas perbandingan kuantitatif dampak integrasi
DFIG-WT dan PMSG-WT terhadap \textit{small signal stability} pada berbagai posisi
bus dan kondisi pembebanan.

\noindent \textbf{Bab V Kesimpulan dan Saran} merangkum temuan penelitian yang
menjawab ketiga rumusan masalah, serta memberikan saran untuk penelitian
selanjutnya terkait integrasi PLTB ke dalam sistem tenaga listrik.

